\chapter*{Introduction: The Organism That Interprets}
\addcontentsline{toc}{chapter}{Introduction: The Organism That Interprets}

An organism that acts must treat some differences as mattering and others as negligible. It must carry expectations about what follows from what, and it must do so before the evidence is complete. Interpretation in this sense is a condition of organized experience, since without a selective ordering of the field there would be no stable objects, no recognizable situations and no actions that could be directed at one outcome rather than another. Yet an ordering that makes life manageable also fixes what can be noticed, and an organization adequate to one range of circumstances can become an impediment when circumstances change or when the organism's own capacities outgrow it. The central problem of this book is how interpretation can be both a condition of organized experience and something that life must repeatedly overcome.

Nietzsche approaches the problem through accounts of schematization, of perspectival knowing and of the multiplicity of drives that coordinate or contest one another beneath the appearance of a unified subject. Predictive processing and active inference approach a related problem from a different direction. They explain perception as the revision of expectations in light of discrepancy, explain action as a means of bringing about the sensory conditions that the organism's model anticipates, and explain learning as the adjustment of the model that generates those expectations \citep{rao1999,friston2010,clark2016,parr2022}. Both bodies of work begin from the embodied activity of a living system and decline to treat cognition as the inspection of a world by a detached observer. That shared starting point is the ground of the comparison. It does not license the claim that either anticipated the other, and the comparison proceeds through shared problems: how simplification is possible without being mere distortion, how a self is constituted as an interpretation of activity, how competing concerns acquire a working hierarchy, and how an organization can be revised without being destroyed.

The governing thesis is stated here provisionally and defended across the chapters. Nietzsche and active inference both make cognition intelligible through the organization of an organism's engagement with its surroundings. Nietzsche additionally asks how that organization acquires authority, and when its preservation becomes an impediment to transformation. Active inference supplies exact vocabulary for the preservation of an organization, namely the maintenance of expected states under a generative model. It is less clear that it supplies an account of the difference between preserving an interpretation that makes life manageable and transforming an interpretation so that new forms of life become possible. Whether that difference can be captured within the formalism, or whether it requires an additional account of value, is the disagreement around which the book is organized.

Four kinds of claim are distinguished throughout, so that resemblance at one level never silently establishes equivalence at another. Historical claims concern what Nietzsche wrote, in which texts, and with what changes of emphasis. Conceptual comparisons concern whether a problem he poses has a counterpart in contemporary modelling and whether the two treatments are compatible, rival or merely parallel. Mathematical reconstructions are proposals, in which a Nietzschean notion is given a precise form that makes its commitments inspectable, and each is accompanied by a statement of what it leaves out. Normative arguments concern which transformations of interpretation deserve affirmation, and these are argued for on their own terms rather than read off either source.

The method has a literary dimension that is more than ornament. Nietzsche's writing is aphoristic, interrupted and addressed to a reader whose expectations it deliberately disturbs. What such writing does to a reader is evidence about the activity of interpretation, and the book attends to it as such, asking how transitions, changes of voice and unresolved ambiguities alter which features of a text a reader treats as authoritative. A generic account of prediction and error could describe some of this, but it could also flatten irony, allusion and productive ambiguity into a single quantity, and the literary chapters are written to keep that risk visible.

The book proceeds in five parts. The first concerns interpretation before the sovereign subject, taking up the textual method, the treatment of knowledge as schematization, the ego as an interpretation of activity, and the drives as competing perspectives. The second supplies the technical foundation of predictive organization and examines action, precision and exploration. The third addresses self-overcoming and the revision of models, including the problem of value creation. The fourth studies literary form as an experiment in interpretation, and the fifth extends the argument to shared models and institutions and marks the limits of the synthesis.
